% Generated from locale/tr/content/model-theory/models-of-arithmetic/introduction.tex; do not hand-edit.


\olfileid[tr]{mod}{mar}{int}
\section{Giriş}

Aritmetiğin \emph{standart modeli}, $\Domain N=\Nat$ olan ve
$\Obj0$, $\prime$, $+$, $\times$ ile $<$ simgelerinin beklendiği gibi
yorumlandığı !!{structure}~$\Struct N$ ifadesidir. Yani $\Obj0$, $0$;
$\prime$, ardıl fonksiyonu; $+$, toplama; $\times$ ise $\Nat$
üzerindeki çarpma olarak yorumlanır. Daha açık olarak:
\begin{align*}
  \Assign{\Obj{0}}{N} & = 0\\
  \Assign{\prime}{N}(n) & = n + 1\\
  \Assign{+}{N}(n, m) & = n + m\\
  \Assign{\times}{N}(n, m) & = nm.
\end{align*}
Elbette $\Lang{L_A}$ için evreni $\Nat$ dışında olan yapılar da
vardır. Örneğin evreni $\Domain M=\{a\}^*$ (tek $a$ simgesinden oluşan
sonlu diziler, yani $\emptyset,a,aa,aaa,\dots$) olan $\Struct M$
yapısını ve şu yorumları alabiliriz:
\begin{align*}
  \Assign{\Obj{0}}{M} & = \emptyset\\
  \Assign{\prime}{M}(s) & = s \concat a\\
  \Assign{+}{M}(n, m) & = a^{n + m}\\
  \Assign{\times}{M}(n, m) & = a^{nm}.
\end{align*}
Bu iki yapı ``özünde aynıdır'': Tek fark !!{domain}s içindeki
!!{element}s ifadeleridir; yorum fonksiyonlarının !!{domain}s
içindeki !!{element}s arasında kurduğu ilişkiler farklı değildir.
İki !!{structure}s ifadesinin \emph{izomorf} olduğunu söyleriz.

Tıkızlık teoreminin kolay bir sonucu olarak, $\Struct N$ içinde doğru
olan her kuramın $\Struct N$ ile izomorf olmayan modelleri de vardır.
Böyle yapılara \emph{standart olmayan} denir. İlginç olan şudur:
Standart bir modelin (yani $\Struct N$ ve onunla izomorf bütün
!!{structure}s ifadelerinin) !!{element}s, standart sayı terimlerinin
$\num n$ değerleriyle tükenir; yani
\[
\Domain N=\Setabs{\Value{\num n}{N}}{n\in\Nat}.
\]
Fakat standart olmayan modellerde durum böyle değildir: $\Struct M$
standart değilse bütün $n$ değerleri için
$x\neq\Value{\num n}{M}$ olan en az bir $x\in\Domain M$ vardır.

Bu standart olmayan elemanlar oldukça ilginçtir: ``sonsuz doğal
sayılardır.'' Varlıkları bir anlamda eksiklik olgularını da açıklar.
Örneğin Peano aritmetiğinin tutarlılık önermesini,
$\OCon[\Th{PA}]$, yani
$\lnot\lexists[x][\OPrf[\Th{PA}](x,\gn{\lfalse})]$ ifadesini ele
alalım. $\Th{PA}$ ne $\OCon[\Th{PA}]$ ne de
$\lnot\OCon[\Th{PA}]$ önermesini ispatladığından ikisi de $\Th{PA}$
kuramına tutarlı biçimde eklenebilir. $\Th{PA}$ tutarlı olduğundan
$\Sat{N}{\OCon[\Th{PA}]}$ ve dolayısıyla
$\Sat/{N}{\lnot\OCon[\Th{PA}]}$. Böylece $\Struct N$,
$\Th{PA}\cup\{\lnot\OCon[\Th{PA}]\}$ kuramının bir modeli değildir
ve bu kuramın bütün modelleri standart olmayan modeller olmak
zorundadır. $\Th{PA}\cup\{\lnot\OCon[\Th{PA}]\}$ modelleri,
$\lexists[x][\OPrf[\Th{PA}](x,\gn{\lfalse})]$ ifadesini doğru kılan
tanık görevindeki bazı !!{element} ifadelerini; yani $\Th{PA}$
kuramından bir çelişkinin !!a{derivation} ifadesinin bir Gödel sayısını
içermelidir. Böyle !!a{element} standart olamaz; çünkü her $n$ için
$\Th{PA}\Proves\lnot\OPrf[\Th{PA}](\num n,\gn{\lfalse})$.

